% GENERATED FILE. Edit canonical lessons, then run npm run generate:books.
% canonical-source-hash: fnv1a64:76582bb08d20c4ff
% canonical-lessons: MW-C165-mistri, MW-C165-repair-note, MW-R165-first-pass-plans-and-notices, MW-R165-second-pass-the-whole-core

\chapter{A Note About a Repair}
\label{ch:mw-a2-a-note-about-a-repair}
\hlchaptermodality{165}

\begin{chapteropening}
\textbf{By the end of this chapter:} \emph{I can read a neighbour's short note about a repair visit and decide when to be at home.}

Complete MW-R165-second-pass-the-whole-core: read the repair note for when to be home, then sweep the whole A2 core once more.
\end{chapteropening}

\section[\texorpdfstring{\mw{मिस्त्री} (mistrī)}{मिस्त्री (mistrī)}]{\mw{मिस्त्री} (mistrī) — a mechanic}

\label{lesson:MW-C165-mistri}



\begin{quote}
Before the new one: say \emph{a neighbour}, \emph{a fan} and \emph{spoiled}. (\textbf{\mw{पाड़ोसी}}, \textbf{\mw{पंखो}}, \textbf{\mw{खराब}}.)
\end{quote}

\subsection*{You'll want to know: \mw{मिस्त्री}}
\textbf{\mw{मिस्त्री}} — \emph{mistrī} — \textquotedblleft{}a mechanic; a mason\textquotedblright{}: the person who comes to mend things.

\begin{quote}
\textbf{\mw{मिस्त्री} \mw{पंखो} \mw{ठीक} \mw{करै} \mw{है।}}
\end{quote}

\begin{quote}
\emph{mistrī paṅkho ṭhīk karai hai.} — The mechanic mends the fan.
\end{quote}

\begin{quote}
\textbf{\mw{मिस्त्री} \mw{काल} \mw{आयो।}}
\end{quote}

\begin{quote}
\emph{mistrī kāl āyo.} — The mechanic came yesterday.
\end{quote}

\begin{cousinweb}
\textbf{\mw{मिस्त्री}} (\emph{mistrī}) is a loan, usually traced to Portuguese \emph{mestre}, \textquotedblleft{}a master (of a trade)\textquotedblright{}.
\end{cousinweb}

\begin{grammarlens}[title={What you've built}]
\textbf{\mw{ठीक} \mw{करणो}}, \emph{to make right}, is how Marwari says \emph{to mend}: \textbf{\mw{ठीक} \mw{करै} \mw{है}}, and in the past with \textbf{\mw{म्हैं}}, \textbf{\mw{ठीक} \mw{कियो}}.
\end{grammarlens}

\subsection*{Your turn}
Say these aloud:

\begin{itemize}
\raggedright
  \item \emph{mistrī}
  \item \textbf{\mw{पंखो} \mw{खराब} \mw{है।}} — \emph{paṅkho kharāb hai.} — The fan is broken.
  \item \textbf{\mw{मिस्त्री} \mw{पंखो} \mw{ठीक} \mw{करै} \mw{है।}} — \emph{mistrī paṅkho ṭhīk karai hai.} — The mechanic mends the fan.
\end{itemize}

\begin{tcolorbox}[breakable,colback=teal!4,colframe=teal!35!black,title={Before you move on}]
What does \textbf{\mw{मिस्त्री}} mean? (\textquotedblleft{}A mechanic; a mason\textquotedblright{}.) Say it once more.
\end{tcolorbox}
\section[(reading)]{(reading) — a note about a repair}

\label{lesson:MW-C165-repair-note}



\begin{quote}
Say \emph{a mechanic} and \emph{a neighbour}.

This one is written to you, by someone you know. Read it for one thing: when you have to be at home.
\end{quote}

\begin{tcolorbox}[breakable,skin=enhanced,colback=orange!6,colframe=orange!45!black,boxrule=0.5pt,arc=1mm,left=6pt,right=6pt,top=4pt,bottom=4pt,fonttitle=\bfseries,title={Read this}]
\begin{quote}
\textbf{\mw{राम}-\mw{राम} \mw{सा}!}
\end{quote}

\begin{quote}
\textbf{\mw{हूं} \mw{पाड़ोसी} \mw{राम} \mw{हूं।}} \textbf{\mw{काल} \mw{सूं} \mw{अठै} \mw{पाणी} \mw{कोनी।}} \textbf{\mw{आज} \mw{मिस्त्री} \mw{कोनी} \mw{आयो।}} \textbf{\mw{हूं} \mw{काल} \mw{दिनूगै} \mw{मिस्त्री} \mw{लावसूं।}} \textbf{\mw{मिस्त्री} \mw{चोखो} \mw{है।}} \textbf{\mw{आप} \mw{काल} \mw{दिनूगै} \mw{९} \mw{सूं} \mw{११} \mw{तांई} \mw{घरै} \mw{रैवो।}} \textbf{\mw{दरवाजो} \mw{खुलो} \mw{राखो।}} \textbf{\mw{दोपार} \mw{हूं} \mw{घरै} \mw{कोनी।}} \textbf{\mw{कोई} \mw{सवाल}? \mw{आज} \mw{आथण} \mw{७} \mw{सूं} \mw{९} \mw{तांई} \mw{हूं} \mw{घरै} \mw{हूं।}} \textbf{\mw{आभार।}} \textbf{\mw{राम}}
\end{quote}

\emph{Read it:} the note once, straight through, without translating

\emph{Read it:} it again, and find the two lines with hours in them
\end{tcolorbox}

\subsection*{What to know first}
\textbf{\mw{रैवो}} and \textbf{\mw{राखो}} are requests, like \textbf{\mw{लावो}}: \emph{please stay}, \emph{please keep}. \textbf{\mw{लावसूं}} is Ram's own future: \emph{I will bring}. \textbf{\mw{काल} \mw{सूं}} means \emph{since yesterday}: \textbf{\mw{सूं}} marks where the time starts.

\subsection*{Your turn}
\begin{itemize}
\raggedright
  \item \emph{Answer:} what is wrong? (There has been no water since yesterday.)
  \item \emph{Answer:} did the mechanic come today? (No: \textbf{\mw{आज} \mw{मिस्त्री} \mw{कोनी} \mw{आयो।}})
  \item \emph{Answer:} if you have a question, when is Ram at home? (This evening, seven until nine.)
\end{itemize}

\begin{tcolorbox}[breakable,colback=teal!4,colframe=teal!35!black,title={Before you move on}]
When must you be at home? (Tomorrow morning, \textbf{\mw{९} \mw{सूं} \mw{११} \mw{तांई}}.) What should you leave open? (The gate: \textbf{\mw{दरवाजो} \mw{खुलो} \mw{राखो।}})
\end{tcolorbox}
\section[(dialogue)]{Plans and notices — a first pass}

\label{lesson:MW-R165-first-pass-plans-and-notices}



\begin{quote}
Say \emph{a mechanic}. Then \emph{I will come tomorrow}.
\end{quote}

\subsection*{Your turn}
\begin{itemize}
\raggedright
  \item \emph{Say it:} \emph{the mechanic came yesterday}; \emph{the mechanic mends the fan}
  \item \emph{Say it:} \emph{perhaps I'll come again tomorrow}; \emph{the day after tomorrow I'll go to the market}
  \item \emph{Say it:} \emph{I won't go tomorrow}; \emph{don't go now}
  \item \emph{Say it:} \emph{the shop is open from six until ten}; \emph{on Monday the shop is closed}
  \item \emph{Say it:} \emph{I heated the water}; \emph{I took milk}
  \item \emph{Answer:} in Ram's note, when must you be at home? (Tomorrow morning, nine until eleven.)
  \item \emph{Read it:} \textbf{\mw{फेर}}, and name its first letter
\end{itemize}

\begin{tcolorbox}[breakable,colback=teal!4,colframe=teal!35!black,title={Before you move on}]
\emph{A mechanic}? (\textbf{\mw{मिस्त्री}}.) \emph{Until}? (\textbf{\mw{तांई}}.) \emph{I will bring}? (\textbf{\mw{लावसूं}}.)
\end{tcolorbox}
\section[(dialogue)]{The whole A2 core — a second pass}

\label{lesson:MW-R165-second-pass-the-whole-core}



\begin{quote}
Say \emph{a mechanic}. When did Ram's note ask you to be at home? (Tomorrow morning.)
\end{quote}

\subsection*{Your turn: people and the present}
Say these aloud:

\begin{itemize}
\raggedright
  \item \emph{he is Ram}; \emph{that is a camel}; \emph{you are here}
  \item \emph{I speak a little}; \emph{Ram reads the newspaper}; \emph{what do you do?}
  \item \emph{I drink milk every day}; \emph{mother goes to the market}
  \item \emph{my cow is here}; \emph{the milk is in the kitchen}; \emph{I come from the shop}
\end{itemize}

\subsection*{Your turn: not, questions, past and future}
\begin{itemize}
\raggedright
  \item \emph{Say it:} \emph{I don't understand}; \emph{Ram isn't here}; \emph{don't go}
  \item \emph{Say it:} \emph{do you recognise this song?}; \emph{yes, I recognise it}; \emph{when do you come?}
  \item \emph{Say it:} \emph{mother went home}; \emph{Ram came yesterday}; \emph{mother got up early}
  \item \emph{Say it:} \emph{I ate a roti}; \emph{I took milk}; \emph{I drew water}
  \item \emph{Say it:} \emph{I will go tomorrow}; \emph{I won't go tomorrow}; \emph{I'll come again}
  \item \emph{Read it:} \textbf{\mw{समझणो}}, \textbf{\mw{ओळखणो}} and \textbf{\mw{काढणो}}, and name the newest letter in each
\end{itemize}

\begin{tcolorbox}[breakable,colback=teal!4,colframe=teal!35!black,title={Before you move on}]
\emph{Went}? (\textbf{\mw{गयो}}, \textbf{\mw{गई}}, \textbf{\mw{गया}}.) \emph{Will go}? (\textbf{\mw{जासूं}}.) \emph{Don't}? (\textbf{\mw{मत}}.) \emph{Not}? (\textbf{\mw{कोनी}}.)
\end{tcolorbox}
